\chapter{2009年安徽卷（文）}

\section{单选题}

\begin{problem}
\(\i\) 是虚数单位，\(\i(1 + \i)\) 等于（\quad）

\choices
    {\(1 + \i\)}
    {\(-1 - \i\)}
    {\(1 - \i\)}
    {\(-1 + \i\)}
\end{problem}

\begin{answer}
D
\end{answer}

\begin{solution}
由 \(\i^2=-1\)，得 \(\i(1+\i)=\i+\i^2=-1+\i\)，故选 D.
\end{solution}

\begin{problem}
若集合 \(A = \{x \mid (2x + 1)(x - 3) < 0\}\)，\(B = \{x \in \N_+ \mid x \leqslant 5\}\)，则 \(A \cap B\) 是（\quad）

\choices
    {\(\{1, 2, 3\}\)}
    {\(\{1, 2\}\)}
    {\(\{4, 5\}\)}
    {\(\{1, 2, 3, 4, 5\}\)}
\end{problem}

\begin{answer}
B
\end{answer}

\begin{solution}
由 \((2x+1)(x-3)<0\) 的两个根为 \(-\frac{1}{2}\) 和 \(3\)，且二次项系数为正，得 \(-\frac{1}{2}<x<3\)，所以 \(A=\left(-\frac{1}{2},3\right)\). 又 \(B=\{1,2,3,4,5\}\)，从而 \(A\cap B=\{1,2\}\)，故选 B.
\end{solution}

\begin{problem}
不等式组 \(
\begin{cases}
x\geqslant0,\\
x+3y\geqslant4,\\
3x+y\leqslant4
\end{cases}
\) 所表示的平面区域的面积等于（\quad）

\choices
    {\( \frac{3}{2} \)}
    {\( \frac{2}{3} \)}
    {\( \frac{4}{3} \)}
    {\( \frac{3}{4} \)}
\end{problem}

\begin{answer}
C
\end{answer}

\begin{solution}
两直线 \(x+3y=4\) 与 \(3x+y=4\) 的交点满足 \(x=y=1\). 当 \(x=0\) 时，两直线分别给出 \(y=\frac{4}{3}\) 和 \(y=4\)，所以所表示区域是顶点为 \((0,\frac{4}{3})\)，\((0,4)\)，\((1,1)\) 的三角形. 其面积为 \(\frac{1}{2}\times\left(4-\frac{4}{3}\right)\times1=\frac{4}{3}\)，故选 C.
\end{solution}

\begin{problem}
“\(a + c > b + d\)”是“\(a > b\)且\(c > d\)”的（\quad）

\choices
    {必要不充分条件}
    {充分不必要条件}
    {充分必要条件}
    {既不充分也不必要条件}
\end{problem}

\begin{answer}
A
\end{answer}

\begin{solution}
记命题 \(P\) 为 \(a+c>b+d\)，记命题 \(Q\) 为 \(a>b\) 且 \(c>d\). 若 \(Q\) 成立，则将两个不等式相加可得 \(a+c>b+d\)，所以 \(P\) 是 \(Q\) 成立的必要条件. 但 \(P\) 不推出 \(Q\)，例如取 \(a=2\)，\(b=1\)，\(c=d=0\)，则 \(P\) 成立而 \(c>d\) 不成立. 因此 \(P\) 是 \(Q\) 的必要不充分条件，故选 A.
\end{solution}

\begin{problem}
已知 \(\{a_{n}\}\) 为等差数列，\(a_{1} + a_{3} + a_{5} = 105\)，\(a_{2} + a_{4} + a_{6} = 99\)，则 \(a_{20}\) 等于（\quad）

\choices
    {\(-1\)}
    {\(1\)}
    {\(3\)}
    {\(7\)}
\end{problem}

\begin{answer}
B
\end{answer}

\begin{solution}
设等差数列的首项为 \(a_1\)，公差为 \(d\). 由题意，\(a_1+a_3+a_5=3(a_1+2d)=105\)，所以 \(a_1+2d=35\)；又 \(a_2+a_4+a_6=3(a_1+3d)=99\)，所以 \(a_1+3d=33\). 两式相减得 \(d=-2\)，代回得 \(a_1=39\). 因而 \(a_{20}=a_1+19d=39-38=1\)，故选 B.
\end{solution}

\begin{problem}
下列曲线中离心率为 \(\frac{\sqrt{6}}{2}\) 的是（\quad）

\choices
    {\(\frac{x^2}{2} - \frac{y^2}{4} = 1\)}
    {\(\frac{x^2}{4} - \frac{y^2}{2} = 1\)}
    {\(\frac{y^2}{4} - \frac{x^2}{6} = 1\)}
    {\(\frac{x^2}{4} -\frac{y^2}{10} = 1\)}
\end{problem}

\begin{answer}
B
\end{answer}

\begin{solution}
双曲线的离心率满足 \(e^2=1+\frac{b^2}{a^2}\)，其中 \(a^2\) 是正平方项的分母，\(b^2\) 是负平方项的分母. 四个选项对应的 \(e^2\) 依次为 \(1+\frac{4}{2}=3\)，\(1+\frac{2}{4}=\frac{3}{2}\)，\(1+\frac{6}{4}=\frac{5}{2}\)，\(1+\frac{10}{4}=\frac{7}{2}\). 由于 \(\left(\frac{\sqrt{6}}{2}\right)^2=\frac{3}{2}\)，故选 B.
\end{solution}

\begin{problem}
直线 \(l\) 过点 \((-1,2)\) 且与直线 \(2x - 3y + 4 = 0\) 垂直，则 \(l\) 的方程是（\quad）

\choices
    {\(3x + 2y - 1 = 0\)}
    {\(3x + 2y + 7 = 0\)}
    {\(2x - 3y + 5 = 0\)}
    {\(2x - 3y + 8 = 0\)}
\end{problem}

\begin{answer}
A
\end{answer}

\begin{solution}
直线 \(2x-3y+4=0\) 的斜率为 \(\frac{2}{3}\)，所以与它垂直的直线 \(l\) 的斜率为 \(-\frac{3}{2}\). 由 \(l\) 经过 \((-1,2)\)，得 \(y-2=-\frac{3}{2}(x+1)\)，化简为 \(3x+2y-1=0\)，故选 A.
\end{solution}

\begin{problem}
设 \(a < b\)，函数 \(y = (x - a)^2 (x - b)\) 的图象可能是（\quad）

\choices
    {\choicebitmap[width=0.15\paperwidth]{img/2009/anhui_liberal/q8_choiceA.png}}
    {\choicebitmap[width=0.15\paperwidth]{img/2009/anhui_liberal/q8_choiceB.png}}
    {\choicebitmap[width=0.15\paperwidth]{img/2009/anhui_liberal/q8_choiceC.png}}
    {\choicebitmap[width=0.15\paperwidth]{img/2009/anhui_liberal/q8_choiceD.png}}
\end{problem}

\begin{answer}
C
\end{answer}

\begin{solution}
函数可写为 \(y=(x-a)^2(x-b)\). 因为 \(a\) 是二重根，图象在 \(x=a\) 处与 \(x\) 轴相切而不穿过；因为 \(b\) 是单根，图象在 \(x=b\) 处穿过 \(x\) 轴. 又函数是首项系数为正的三次函数，所以当 \(x\to-\infty\) 时 \(y\to-\infty\)，当 \(x\to+\infty\) 时 \(y\to+\infty\). 四个图象中只有 C 同时满足这些特征，故选 C.
\end{solution}

\begin{problem}
设函数 \( f(x) = \frac{\sin \theta}{3} x^3 + \frac{\sqrt{3} \cos \theta}{2} x^2 + \tan \theta \)，其中 \( \theta \in \left[0, \frac{5\pi}{12}\right] \)，则导数 \( f'(1) \) 的取值范围是（\quad）

\choices
    {\([-2, 2]\)}
    {\([\sqrt{2},\sqrt{3}]\)}
    {\([\sqrt{3}, 2]\)}
    {\([\sqrt{2}, 2]\)}
\end{problem}

\begin{answer}
D
\end{answer}

\begin{solution}
对 \(x\) 求导得 \(f'(x)=\sin\theta\,x^2+\sqrt{3}\cos\theta\,x\)，所以 \(f'(1)=\sin\theta+\sqrt{3}\cos\theta=2\sin\left(\theta+\frac{\pi}{3}\right)\). 令 \(\varphi=\theta+\frac{\pi}{3}\)，则 \(\varphi\in\left[\frac{\pi}{3},\frac{3\pi}{4}\right]\). 在该区间上，\(\sin\varphi\) 的最大值为 \(1\)，最小值为 \(\sin\frac{3\pi}{4}=\frac{\sqrt{2}}{2}\)，故 \(f'(1)\in[\sqrt{2},2]\)，选 D.
\end{solution}

\begin{problem}
考察正方体 \(6\) 个面的中心，从中任意选 \(3\) 个点连成三角形，再把剩下的 \(3\) 个点也连成三角形，则所得的两个三角形全等的概率等于（\quad）

\choices
    {\(1\)}
    {\( \frac{1}{2} \)}
    {\( \frac{1}{3} \)}
    {\(0\)}
\end{problem}

\begin{answer}
A
\end{answer}

\begin{solution}
把六个面中心分别记为 \((\pm1,0,0)\)，\((0,\pm1,0)\)，\((0,0,\pm1)\). 任取三个面中心有两种情况：若三个点分别来自三组相对面，则任意两点之间的距离都为 \(\sqrt{2}\)，所得三角形为等边三角形；剩下的三个点也分别来自三组相对面，同样组成等边三角形. 若三个点中包含一对相对面中心，则第三个点来自另一组相对面，所得三角形的三边为 \(2\)，\(\sqrt{2}\)，\(\sqrt{2}\)；剩下的三个点也包含一对相对面中心，且三边同样为 \(2\)，\(\sqrt{2}\)，\(\sqrt{2}\). 因而无论哪种取法，两个三角形都全等，事件必然发生，概率为 \(1\)，故选 A.
\end{solution}

\section{填空题}

\begin{problem}
在空间直角坐标系中，已知点 \(A(1,0,2)\)，\(B(1,-3,1)\)，点 \(M\) 在 \(y\) 轴上，且 \(M\) 到 \(A\) 与到 \(B\) 的距离相等，则 \(M\) 的坐标是\fillinblank{}.
\end{problem}

\begin{answer}
\((0,-1,0)\)
\end{answer}

\begin{solution}
设 \(M=(0,t,0)\)，则 \(MA^2=(0-1)^2+(t-0)^2+(0-2)^2=t^2+5\)，\(MB^2=(0-1)^2+(t+3)^2+(0-1)^2=t^2+6t+11\). 由 \(MA=MB\)，得 \(t^2+5=t^2+6t+11\)，所以 \(t=-1\). 因此 \(M=(0,-1,0)\).
\end{solution}

\begin{problem}
程序框图（即算法流程图）如图所示，其输出结果是\fillinblank{}.

\bitmapfigure[width=4.5cm]{img/2009/anhui_liberal/q12_fig1.png}
\end{problem}

\begin{answer}
127
\end{answer}

\begin{solution}
由流程图，初值为 \(a_0=1\)，每循环一次先将 \(a\) 更新为 \(2a+1\)，再判断是否满足 \(a>100\). 记更新 \(k\) 次后的数为 \(a_k\)，则 \(a_{k+1}=2a_k+1\). 因而 \(a_{k+1}+1=2(a_k+1)\)，从而 \(a_k+1=2^k(a_0+1)=2^{k+1}\)，即 \(a_k=2^{k+1}-1\). 计算得 \(a_5=63\leqslant100\)，而 \(a_6=127>100\)，所以流程图输出 \(127\).
\end{solution}

\begin{problem}
从长度分别为 \(2\)，\(3\)，\(4\)，\(5\) 的四条线段中任意取出三条，则以这三条线段为边可以构成三角形的概率是\fillinblank{}.
\end{problem}

\begin{answer}
\(\frac{3}{4}\)
\end{answer}

\begin{solution}
从四条线段中任取三条共有 \(\mathrm C_4^3=4\) 种取法. 三条线段能构成三角形的条件是最长边小于另外两边之和. 四种取法 \(\{2,3,4\}\)，\(\{2,3,5\}\)，\(\{2,4,5\}\)，\(\{3,4,5\}\) 中，只有 \(2+3=5\) 的取法不能构成三角形，其余三种均能构成三角形. 因此所求概率为 \(\frac{3}{4}\).
\end{solution}

\begin{problem}
在平行四边形 \(ABCD\) 中，\(E\) 和 \(F\) 分别是边 \(CD\) 和 \(BC\) 的中点. 若 \(\overrightarrow{AC} = \lambda \overrightarrow{AE} + \mu \overrightarrow{AF}\)，其中 \(\lambda\)，\(\mu \in \R\)，则 \(\lambda + \mu =\)\fillinblank{}.
\end{problem}

\begin{answer}
\(\frac{4}{3}\)
\end{answer}

\begin{solution}
以 \(A\) 为原点，令 \(\overrightarrow{AB}=\bs{b}\)，\(\overrightarrow{AD}=\bs{d}\)，则 \(\overrightarrow{AC}=\bs{b}+\bs{d}\). 因为 \(E\)，\(F\) 分别是 \(CD\)，\(BC\) 的中点，所以 \(\overrightarrow{AE}=\frac{1}{2}\bs{b}+\bs{d}\)，\(\overrightarrow{AF}=\bs{b}+\frac{1}{2}\bs{d}\). 代入题设得
\[
\bs{b}+\bs{d}=\lambda\left(\frac{1}{2}\bs{b}+\bs{d}\right)+\mu\left(\bs{b}+\frac{1}{2}\bs{d}\right).
\]
比较 \(\bs{b}\)，\(\bs{d}\) 的系数，得到 \(\frac{\lambda}{2}+\mu=1\)，\(\lambda+\frac{\mu}{2}=1\). 两式相减得 \(\lambda=\mu\)，所以 \(\lambda=\mu=\frac{2}{3}\)，从而 \(\lambda+\mu=\frac{4}{3}\).
\end{solution}

\begin{problem}
对于四面体 \(ABCD\)，下列命题正确的是\fillinblank{}.（写出所有正确命题的编号）

\begin{circlelist}
\item 相对棱 \(AB\) 与 \(CD\) 所在的直线异面；
\item 由顶点 \(A\) 作四面体的高，其垂足是 \(\triangle BCD\) 的三条高线的交点；
\item 若分别作 \(\triangle ABC\) 和 \(\triangle ABD\) 的边 \(AB\) 上的高，则这两条高所在直线异面；
\item 任何三个面的面积之和都大于第四个面的面积；
\item 分别作三组相对棱中点的连线，所得的三条线段相交于一点.
\end{circlelist}
\end{problem}

\begin{answer}
\(\circled{1}\)，\(\circled{4}\)，\(\circled{5}\)
\end{answer}

\begin{solution}
命题\(\circled{1}\)正确：四面体的两条相对棱不共面，所以它们所在的直线异面.

命题\(\circled{2}\)错误. 顶点 \(A\) 到平面 \(BCD\) 的垂足只是 \(A\) 在平面 \(BCD\) 上的射影，一般不必是三角形 \(BCD\) 的垂心. 例如取 \(B=(0,0,0)\)，\(C=(4,0,0)\)，\(D=(0,3,0)\)，\(A=(1,1,1)\)，则三角形 \(BCD\) 的垂心为 \(B\)，而 \(A\) 到平面 \(BCD\) 的垂足为 \((1,1,0)\)，二者不重合.

命题\(\circled{3}\)错误. 例如取 \(A=(0,0,0)\)，\(B=(1,0,0)\)，\(C=(0,1,0)\)，\(D=(0,0,1)\). 此时三角形 \(ABC\) 和 \(ABD\) 中，边 \(AB\) 上的高都以 \(A\) 为垂足，两条高所在直线分别为 \(AC\) 和 \(AD\)，它们相交于 \(A\)，不是异面直线.

命题\(\circled{4}\)正确. 设四个面的面积向量分别为 \(\bs{S}_1\)，\(\bs{S}_2\)，\(\bs{S}_3\)，\(\bs{S}_4\)，取向均为四面体的外法向，则有 \(\bs{S}_1+\bs{S}_2+\bs{S}_3+\bs{S}_4=\bs{0}\). 因而 \(\bs{S}_4=-(\bs{S}_1+\bs{S}_2+\bs{S}_3)\)，由三角不等式得 \(S_4\leqslant S_1+S_2+S_3\). 非退化四面体的三个相邻面外法向不共线，三角不等式取严格不等式，故任意一个面的面积都小于另外三个面的面积之和.

命题\(\circled{5}\)正确. 设四个顶点的位置向量为 \(\bs{A}\)，\(\bs{B}\)，\(\bs{C}\)，\(\bs{D}\). 连接 \(AB\)，\(CD\) 的中点时，所得线段的中点为 \(\frac{1}{4}(\bs{A}+\bs{B}+\bs{C}+\bs{D})\)；连接 \(AC\)，\(BD\) 或 \(AD\)，\(BC\) 的中点时，所得线段的中点也都是该点. 所以三条线段相交于一点. 因此正确命题的编号为 \(\circled{1}\)，\(\circled{4}\)，\(\circled{5}\).
\end{solution}

\section{解答题}

\begin{problem}
在 \(\triangle ABC\) 中，\(C - A = \frac{\pi}{2}\)，\(\sin B = \frac{1}{3}\).
\begin{enumerate}
\item 求 \( \sin A \) 的值；
\item 设 \( AC = \sqrt{6} \)，求 \( \triangle ABC \) 的面积.
\end{enumerate}
\end{problem}

\begin{solution}
\begin{enumerate}
\item 由三角形内角和以及 \(C-A=\frac{\pi}{2}\)，得
\[
B=\pi-A-C=\frac{\pi}{2}-2A.
\]
因为 \(A>0\)，\(B>0\)，所以 \(0<A<\frac{\pi}{4}\). 于是
\[
\sin B=\sin\left(\frac{\pi}{2}-2A\right)=\cos 2A
=1-2\sin^2 A=\frac{1}{3}.
\]
故 \(\sin^2 A=\frac{1}{3}\). 又 \(A\) 为三角形内角，\(\sin A>0\)，所以
\[
\sin A=\frac{\sqrt{3}}{3}.
\]
\item 设 \(a=BC\)，\(b=CA=\sqrt{6}\)，\(c=AB\). 由正弦定理，
\[
\frac{b}{\sin B}=\frac{\sqrt{6}}{1/3}=3\sqrt{6}.
\]
又
\[
\sin C=\sin\left(A+\frac{\pi}{2}\right)=\cos A
=\sqrt{1-\sin^2 A}=\frac{\sqrt{6}}{3},
\]
所以 \(c=3\sqrt{6}\cdot\frac{\sqrt{6}}{3}=6\). 因此
\[
S_{\triangle ABC}=\frac{1}{2}\,AC\cdot AB\cdot\sin A
=\frac{1}{2}\cdot\sqrt{6}\cdot6\cdot\frac{\sqrt{3}}{3}
=3\sqrt{2}.
\]
\end{enumerate}
\end{solution}

\begin{problem}
某良种培育基地正在培育一种小麦新品种 A，将其与原有的一个优良品种 B 进行对照试验，两种小麦各种植了 \(25\) 亩，所得亩产数据（单位：千克）如下：

\begin{center}
\begin{tabular}{c|c}
品种 A（千克） & 品种 B（千克） \\
\hline
\(357\) & \(363\) \\
\(359\) & \(371\) \\
\(367\) & \(374\) \\
\(368\) & \(383\) \\
\(375\) & \(385\) \\
\(388\) & \(386\) \\
\(392\) & \(391\) \\
\(399\) & \(392\) \\
\(400\) & \(394\) \\
\(405\) & \(395\) \\
\(412\) & \(397\) \\
\(414\) & \(397\) \\
\(415\) & \(400\) \\
\(421\) & \(401\) \\
\(423\) & \(401\) \\
\(423\) & \(403\) \\
\(427\) & \(406\) \\
\(430\) & \(407\) \\
\(430\) & \(410\) \\
\(434\) & \(412\) \\
\(443\) & \(415\) \\
\(445\) & \(416\) \\
\(451\) & \(422\) \\
\(454\) & \(430\) \\
\end{tabular}
\end{center}

\begin{enumerate}
\item 完成所附的茎叶图；
\item 用茎叶图处理现有的数据，有什么优点？
\item 通过观察茎叶图，对品种 A 与 B 的亩产量及其稳定性进行比较，写出统计结论.
\end{enumerate}

\begin{center}
\begin{tabular}{p{3.5cm}|c|p{3.5cm}}
\multicolumn{1}{c|}{A} & & \multicolumn{1}{c}{B} \\
\hline
 & \(35\) & \\
 & \(36\) & \\
 & \(37\) & \\
 & \(38\) & \\
 & \(39\) & \\
 & \(40\) & \\
 & \(41\) & \\
 & \(42\) & \\
 & \(43\) & \\
 & \(44\) & \\
 & \(45\) & \\
\end{tabular}
\end{center}
\end{problem}

\begin{solution}
\begin{enumerate}
\item 以数据的百位和十位组成茎，以个位作为叶. 例如 \(357\) 记作茎 \(35\)、叶 \(7\)，\(359\) 记作茎 \(35\)、叶 \(9\). 对全部数据按茎从小到大排列，并将同一茎上的叶按从小到大排列，得到完整茎叶图：

\begin{center}
\begin{tabular}{p{3.5cm}|c|p{3.5cm}}
\multicolumn{1}{c|}{A} & & \multicolumn{1}{c}{B} \\
\hline
\(7\enspace9\) & \(35\) & \\
\(7\enspace8\) & \(36\) & \(3\) \\
\(5\) & \(37\) & \(1\enspace4\) \\
\(8\) & \(38\) & \(3\enspace5\enspace6\) \\
\(2\enspace9\) & \(39\) & \(1\enspace2\enspace4\enspace5\enspace7\enspace7\) \\
\(0\enspace5\) & \(40\) & \(0\enspace1\enspace1\enspace3\enspace6\enspace7\) \\
\(2\enspace4\enspace5\) & \(41\) & \(0\enspace2\enspace5\enspace6\) \\
\(1\enspace3\enspace3\enspace7\) & \(42\) & \(2\) \\
\(0\enspace0\enspace4\) & \(43\) & \(0\) \\
\(3\enspace5\) & \(44\) & \\
\(1\enspace4\) & \(45\) & \\
\end{tabular}
\end{center}
\item 茎叶图既能保留原始数据的全部信息，又能直观表示数据的集中趋势和离散程度，便于比较两组数据的大小、分布及稳定性.
\item 从图中看，品种 A 的数据整体更偏向较大的茎，亩产量总体高于品种 B；品种 A 的亩产量范围为 \(357\) 至 \(454\)，极差为 \(454-357=97\)，品种 B 的亩产量范围为 \(363\) 至 \(430\)，极差为 \(430-363=67\). 因此，品种 A 的亩产量总体较高，但波动较大；品种 B 的亩产量总体较低，但稳定性较好.
\end{enumerate}
\end{solution}

\begin{problem}
已知椭圆 \(\frac{x^2}{a^2} + \frac{y^2}{b^2} = 1 \) （\(a > b > 0\)）的离心率为 \(\frac{\sqrt{3}}{3}\)，以原点为圆心，椭圆短半轴长为半径的圆与直线 \(y = x + 2\) 相切.
\begin{enumerate}
\item 求 \(a\) 与 \(b\).
\item 设该椭圆的左，右焦点分别为 \(F_{1}\) 和 \(F_{2}\)，直线 \(l_{1}\) 过 \(F_{2}\) 且与 \(x\) 轴垂直，动直线 \(l_{2}\) 与 \(y\) 轴垂直，交 \(l_{1}\) 于点 \(P\). 求线段 \(PF_{1}\) 的垂直平分线与 \(l_{2}\) 的交点 \(M\) 的轨迹方程，并指明曲线类型.
\end{enumerate}
\end{problem}

\begin{solution}
\begin{enumerate}
\item 由题意，圆的半径为 \(b\). 直线 \(y=x+2\) 化为 \(x-y+2=0\)，原点到该直线的距离为
\[
\frac{2}{\sqrt{1^2+(-1)^2}}=\sqrt{2}.
\]
因为圆与直线相切，所以 \(b=\sqrt{2}\). 设椭圆的半焦距为 \(c\)，则
\[
\frac{c}{a}=\frac{\sqrt{3}}{3},\qquad c^2=a^2-b^2.
\]
于是
\[
\frac{a^2-b^2}{a^2}=\frac{1}{3},\qquad
\frac{b^2}{a^2}=\frac{2}{3}.
\]
代入 \(b^2=2\)，得 \(a^2=3\)，所以
\[
a=\sqrt{3},\qquad b=\sqrt{2}.
\]
\item 由上问 \(c=\sqrt{a^2-b^2}=1\)，故 \(F_{1}(-1,0)\)，\(F_{2}(1,0)\). 直线 \(l_{1}\) 的方程为 \(x=1\). 设动直线 \(l_{2}\) 为 \(y=t\)，则
\[
P=(1,t).
\]
设交点 \(M=(x,t)\). 因为 \(M\) 在 \(PF_{1}\) 的垂直平分线上，所以 \(MP=MF_{1}\)，从而
\[
(x-1)^2=(x+1)^2+t^2.
\]
化简得
\[
-4x=t^2,\qquad x=-\frac{t^2}{4}.
\]
又 \(M\) 的坐标满足 \(y=t\)，消去参数 \(t\)，得到
\[
y^2=-4x.
\]
因此 \(M\) 的轨迹是顶点在原点、开口向左、焦点在 \(x\) 轴上的抛物线.
\end{enumerate}
\end{solution}


\begin{problem}
已知数列 \(\{a_n\}\) 的前 \(n\) 项和 \(S_{n} = 2n^{2} + 2n\)，数列 \(\{b_n\}\) 的前 \(n\) 项和 \(T_{n} = 2 - b_{n}\).
\begin{enumerate}
\item 求数列 \( \{a_{n}\} \) 与 \( \{b_{n}\} \) 的通项公式；
\item 设 \(c_{n} = a_{n}^{2} \cdot b_{n}\)，证明：当且仅当 \(n \geqslant 3\) 时，\(c_{n+1} < c_{n}\).
\end{enumerate}
\end{problem}

\begin{solution}
\begin{enumerate}
\item 当 \(n=1\) 时，\(a_1=S_1=2+2=4\). 当 \(n\geqslant2\) 时，
\[
a_n=S_n-S_{n-1}
=\left(2n^2+2n\right)-\left(2(n-1)^2+2(n-1)\right)
=4n.
\]
所以对一切 \(n\in\N_+\)，都有 \(a_n=4n\).

由 \(T_1=b_1=2-b_1\)，得 \(b_1=1\). 当 \(n\geqslant2\) 时，
\[
b_n=T_n-T_{n-1}
=\left(2-b_n\right)-\left(2-b_{n-1}\right)
=b_{n-1}-b_n,
\]
所以 \(2b_n=b_{n-1}\)，即 \(b_n=\frac{1}{2}b_{n-1}\). 递推可得
\[
b_n=\frac{1}{2^{n-1}}=2^{1-n}.
\]
因此
\[
a_n=4n,\qquad b_n=2^{1-n}.
\]
\item 由上问，
\[
c_n=a_n^2b_n=(4n)^2\cdot2^{1-n}.
\]
因为 \(c_n>0\)，所以
\[
\frac{c_{n+1}}{c_n}
=\frac{16(n+1)^2\cdot2^{-n}}{16n^2\cdot2^{1-n}}
=\frac{(n+1)^2}{2n^2}.
\]
于是
\[
c_{n+1}<c_n
\Longleftrightarrow
\frac{(n+1)^2}{2n^2}<1
\Longleftrightarrow
(n+1)^2<2n^2
\Longleftrightarrow
n^2-2n-1>0.
\]
由于 \(n\in\N_+\)，上式等价于 \(n>1+\sqrt{2}\)，也等价于 \(n\geqslant3\). 当 \(n=1,2\) 时不等式不成立，故当且仅当 \(n\geqslant3\) 时，\(c_{n+1}<c_n\).
\end{enumerate}
\end{solution}

\begin{problem}
如图，\(ABCD\) 的边长为 \(2\) 的正方形，直线 \(l\) 与平面 \(ABCD\) 平行，\(E\) 和 \(F\) 是 \(l\) 上的两个不同点，且 \(EA = ED\)，\(FB = FC\). \(E'\) 和 \(F'\) 是平面 \(ABCD\) 内的两点，\(EE'\) 和 \(FF'\) 都与平面 \(ABCD\) 垂直.

\begin{enumerate}
\item 证明：直线 \(E'F'\) 垂直且平分线段 \(AD\)；
\item 若 \(\angle EAD = \angle EAB = 60^{\circ}\)，\(EF = 2\)，求多面体 \(ABCDEF\) 的体积.
\end{enumerate}

\bitmapfigure[max width=0.42\linewidth]{img/2009/anhui_liberal/q19_fig1.png}
\end{problem}

\begin{solution}
\begin{enumerate}
\item 建立空间直角坐标系，令平面 \(ABCD\) 为 \(z=0\)，并取
\[
A=(0,0,0),\quad B=(2,0,0),\quad D=(0,2,0),\quad C=(2,2,0).
\]
设 \(E'=(u,v,0)\). 由 \(EA=ED\)，且 \(EE'\perp\) 平面 \(ABCD\)，得 \(E'A=E'D\)，所以
\[
u^2+v^2=u^2+(v-2)^2,
\]
从而 \(v=1\). 同理，设 \(F'=(r,s,0)\)，由 \(FB=FC\) 得 \(F'B=F'C\)，于是
\[
(r-2)^2+s^2=(r-2)^2+(s-2)^2,
\]
从而 \(s=1\). 因此 \(E'\)，\(F'\) 都在平面内的直线 \(y=1\) 上. 由于 \(E\)，\(F\) 是直线 \(l\) 上的两个不同点，且 \(l\) 平行于平面 \(ABCD\)，故 \(E'\ne F'\)，所以直线 \(E'F'\) 就是 \(y=1\). 它经过 \(AD\) 的中点 \((0,1,0)\)，且 \(AD\) 的方向为 \(y\) 轴方向，故 \(E'F'\perp AD\)，并且平分 \(AD\).
\item 设 \(EE'=h\). 因为 \(l\parallel\) 平面 \(ABCD\)，所以 \(E\) 与 \(F\) 到平面的距离相等. 由上问，\(E'\) 的坐标可写为 \((u,1,0)\)，于是 \(E=(u,1,h)\)，且
\[
AE=\sqrt{u^2+1+h^2}.
\]
由 \(\angle EAD=60^\circ\)，有
\[
\cos\angle EAD
=\frac{\overrightarrow{AE}\cdot\overrightarrow{AD}}{AE\cdot AD}
=\frac{1}{AE}=\frac{1}{2},
\]
所以 \(AE=2\). 再由 \(\angle EAB=60^\circ\)，得
\[
\cos\angle EAB
=\frac{\overrightarrow{AE}\cdot\overrightarrow{AB}}{AE\cdot AB}
=\frac{u}{AE}=\frac{1}{2},
\]
从而 \(u=1\). 代入 \(AE^2=u^2+1+h^2\)，得
\[
4=1+1+h^2,\qquad h=\sqrt{2}.
\]
由第一问，\(E'F'\parallel AB\)，而 \(EF\) 是其在平行平面上的对应线段，所以 \(EF\parallel AB\). 按图中 \(E,F\) 的取向，\(\overrightarrow{EF}=\overrightarrow{AB}\)；又 \(EF=AB=2\)，因此 \(F\) 是将 \(E\) 沿 \(AB\) 方向平移 \(2\) 个单位得到的点，多面体 \(ABCDEF\) 是以 \(\triangle ADE\) 为底、以 \(EF\) 为棱的三棱柱. 由上面得到 \(AE=2\)，所以
\[
S_{\triangle ADE}=\frac{1}{2}\cdot AD\cdot AE\cdot\sin\angle EAD
=\frac{1}{2}\cdot2\cdot2\cdot\sin60^{\circ}=\sqrt{3}.
\]
故
\[
V_{ABCDEF}=S_{\triangle ADE}\cdot EF=\sqrt{3}\cdot2=2\sqrt{3}.
\]
\end{enumerate}
\end{solution}

\begin{problem}
已知函数 \(f(x) = x - \frac{2}{x} + 1 - a \ln x\)，\(a > 0\).
\begin{enumerate}
\item 讨论 \( f(x) \) 的单调性；
\item 设 \(a = 3\)，求 \(f(x)\) 在区间 \([1, \e^2]\) 上值域，其中 \(\e = 2.71828 \cdots\) 是自然对数的底数.
\end{enumerate}
\end{problem}

\begin{solution}
\begin{enumerate}
\item 函数的定义域为 \(x>0\)，且
\[
f'(x)=1+\frac{2}{x^2}-\frac{a}{x}
=\frac{x^2-ax+2}{x^2}.
\]
因为 \(x^2>0\)，所以 \(f'(x)\) 的符号由 \(g(x)=x^2-ax+2\) 决定.

当 \(0<a<2\sqrt{2}\) 时，\(g(x)\) 的判别式 \(a^2-8<0\)，且二次项系数为正，所以 \(g(x)>0\). 因此 \(f(x)\) 在 \((0,+\infty)\) 上单调递增.

当 \(a=2\sqrt{2}\) 时，\(g(x)=(x-\sqrt{2})^2\geqslant0\)，所以 \(f(x)\) 在 \((0,+\infty)\) 上单调递增.

当 \(a>2\sqrt{2}\) 时，令
\[
x_1=\frac{a-\sqrt{a^2-8}}{2},\qquad
x_2=\frac{a+\sqrt{a^2-8}}{2},
\]
则 \(0<x_1<x_2\)，并且
\[
g(x)=(x-x_1)(x-x_2).
\]
所以 \(f(x)\) 在 \((0,x_1)\) 和 \((x_2,+\infty)\) 上单调递增，在 \((x_1,x_2)\) 上单调递减.
\item 当 \(a=3\) 时，
\[
f'(x)=\frac{x^2-3x+2}{x^2}
=\frac{(x-1)(x-2)}{x^2}.
\]
因此 \(f(x)\) 在 \([1,2]\) 上单调递减，在 \([2,+\infty)\) 上单调递增. 考察区间 \([1,\e^2]\) 上的端点和极值点：
\[
f(1)=0,\qquad
f(2)=2-3\ln2,\qquad
f(\e^2)=\e^2-\frac{2}{\e^2}-5.
\]
其中 \(2-3\ln2\approx-0.0794\)，\(\e^2-\frac{2}{\e^2}-5\approx2.118\)，故最小值为 \(2-3\ln2\)，最大值为 \(\e^2-\frac{2}{\e^2}-5\). 所以值域为
\[
\left[2-3\ln2,\e^2-\frac{2}{\e^2}-5\right].
\]
\end{enumerate}
\end{solution}
